\section{Sequential collection and stopping: proofs}
\label{sec:annex-p15-sequential-evidence}

\begin{annexitem}{annex:p15.se.posterior}{thm:p15-se-posterior}{Odds, thresholds and exclusive decisions}
\begin{proof}
The two joint prior-weighted densities of the history are
\(g_0p_{0m}\) and \(g_1p_{1m}\). Divide each by their sum to
obtain the posterior probabilities. Their ratio is
\((g_1/g_0)L_m\). The inequality
\(g_1p_{1m}/(g_0p_{0m}+g_1p_{1m})\geq d_1\) rearranges to
the upper ratio boundary; the analogous null inequality gives the
lower boundary. Both cannot hold because their posterior probabilities
sum to one while \(d_0+d_1>1\). Histories with zero total density
are null in the prior mixture and can be assigned any decision.
This derivation establishes posterior conditions only; the two fixed-law
sampling error probabilities are not those conditional probabilities.
\end{proof}
\end{annexitem}

\begin{annexitem}{annex:p15.se.termination}{thm:p15-se-termination}{Affinity tails and a block argument under a third law}
\begin{proof}
Cauchy--Schwarz gives \(\eta_0\leq1\). Equality requires the
square-root densities to be proportional almost everywhere. Their unit
integrals force equality of the densities and hence of the laws, contrary
to the hypothesis. Thus \(\eta_0<1\). Under the null, with singular
values assigned on null sets,
\(E_0\sqrt{f_1/f_0}=\eta_0\); under the alternative,
\(E_1\sqrt{f_0/f_1}=\eta_0\). Independence gives the corresponding
\(m\)-observation expectations \(\eta_0^m\). On \(\{T>m\}\),
the history has no singular ratio and \(B<L_m<A\). Markov's inequality
therefore yields both stated survival bounds. If \(\eta_0=0\),
the first observation has disjoint support and the test stops immediately.
Otherwise Tonelli's identity
\(E_iT=\sum_{m\geq0}P_i(T>m)\) and the geometric bounds give a finite
mean. They also give \(P_i(T=\infty)=0\). Coincident laws have ratio
one on every finite history almost surely. They never leave an interval
containing one; otherwise the first observation is already terminal.

For the third-law assertion, some \(\varepsilon>0\) has either
\(P(z_1\geq\varepsilon)=p>0\) or
\(P(z_1\leq-\varepsilon)=p>0\); otherwise the score would be zero
almost surely. Choose an integer \(k\) with
\(k\varepsilon>a-b\). From any score sum in the continuation interval,
a block of \(k\) increments all satisfying the chosen inequality
forces exit. That block has probability \(p^k\), independently of
the past. Starting after the first observation handles even an initial
zero sum outside the interval: conditional on survival it is inside.
Hence \(P(T>1+jk)\leq(1-p^k)^j\). Summing this geometric
tail proves finite mean without a moment assumption on the increments.
An identically zero score cannot exit when \(0\in(b,a)\);
otherwise its first observation is already terminal.
\end{proof}
\end{annexitem}

\begin{annexitem}{annex:p15.se.errors}{thm:p15-se-errors}{The full stopped-history density comparison}
\begin{proof}
Let \(E_m^i\) be the measurable region of histories that first stop
at stage \(m\) with decision \(i\). Previous continuation conditions
are part of this region. These stopping events are disjoint over
stages. On \(E_m^1\), \(p_{1m}\geq Ap_{0m}\); on
\(E_m^0\), \(p_{1m}\leq Bp_{0m}\), including one-sided zeros.
Integrating against \(\nu^{\otimes m}\) and summing nonnegative
quantities gives
\(P_1(Q_1)\geq AP_0(Q_1)\) and
\(P_1(Q_0)\leq BP_0(Q_0)\). The established termination under both
laws identifies these four probabilities with
\(1-\beta',\alpha',\beta',1-\alpha'\), respectively.
This proves the coupled inequalities. Dropping the factors no larger
than one gives the separate bounds. The two bounding lines meet at
\(\alpha'=(1-B)/(A-B)\),
\(\beta'=B(A-1)/(A-B)\); their nonnegative-axis intersections give
the other three vertices. This proves the stated quadrilateral.
The proposed conservative boundaries
then imply both target inequalities. No infinite stopped-mean identity
was used, and no unobserved future entered the stopping regions.
\end{proof}
\end{annexitem}

\begin{annexitem}{annex:p15.se.dependent}{thm:p15-se-dependent}{Only the correct finite-history densities are required}
\begin{proof}
Repeat the preceding integration on \(E_m^i\) using the actual
joint densities of the two process laws. Each ratio inequality is
pointwise on the same observed-history region. Disjointness of stopping
events and almost-sure termination give the same sum identities.
No factorization or independence occurs in this proof. Consistent
marginals ensure that the finite-history densities belong to the
specified process laws. A product of dependent marginal densities is
a different joint law and cannot be substituted in these integrals.
\end{proof}
\end{annexitem}

\begin{annexitem}{annex:p15.se.nominal}{thm:p15-se-nominal}{Both coupled bounds and the explicit one-step failure}
\begin{proof}
Substitute the nominal boundaries into the exact coupled inequalities:
\((1-\beta)\alpha'\leq\alpha(1-\beta')\) and
\((1-\alpha)\beta'\leq\beta(1-\alpha')\).
The separate upper bounds follow by dropping the actual complements.
Adding the two inequalities cancels their cross terms and gives
\(\alpha'+\beta'\leq\alpha+\beta\). This controls the sum, not each
target separately. In the stated Bernoulli example a success has ratio
\((99/100)/(11/100)=9\); a failure has ratio
\((1/100)/(89/100)=1/89<1/9\). Both outcomes stop at stage one.
The upper decision has null probability \(11/100\), and the lower
decision has alternative probability \(1/100\). All assertions follow
by direct calculation, including the equality-at-boundary convention.
\end{proof}
\end{annexitem}

\begin{annexitem}{annex:p15.se.mean}{thm:p15-se-mean}{Predictability, absolute convergence and conditional excess bounds}
\begin{proof}
The event \(\{T\geq j\}\) depends only on information before the
\(j\)-th observation and is independent of its score. Thus
\[
 E\sum_{j\geq1}|z_j|{\bf1}_{\{T\geq j\}}
     =E|z_1|\sum_{j\geq1}P(T\geq j)=E|z_1|ET<\infty.
\]
Absolute integration permits interchange for the signed stopped sum,
giving \(EZ_T=\mu ET\). This also proves terminal integrability.
Partition by the two terminal decisions to get the exact decomposition.

At a history still continuing before stage \(j\), an upper exit
requires \(z_j\geq a-Z_{j-1}=u>0\). Its excess over the upper boundary
is \(z_j-u\). Conditional on that history and on the exit at that
stage, its conditional mean is no larger than \(\xi\). Independence
justifies use of the same one-step score distribution for every history.
Mixing these inequalities over stages and histories gives
\(c_+\leq\xi\). A lower exit requires
\(z_j\leq-(Z_{j-1}-b)=-u\), and its signed excess
\(Z_T-b=z_j+u\) has conditional mean at least \(\xi'\).
Hence \(-c_-\geq\xi'\). Substitution proves the two inequalities.
Division by a negative mean reverses them, and division by zero is not
allowed. This is also the integrable-increment form of
\cref{thm:p15-os-wald-first}.
\end{proof}
\end{annexitem}

\begin{annexitem}{annex:p15.se.zero-drift}{thm:p15-se-zero-drift}{Second moments and both zero-overshoot formulas}
\begin{proof}
Apply \cref{thm:p15-os-wald-second} to the centered independent
scores and the finite expected stopping time. It establishes terminal
square integrability and \(EZ_T^2=\sigma_z^2ET\), so partitioning
by the terminal event gives the displayed identity with its finite
conditional terms. Under zero overshoot, the first stopped identity is
\(0=(1-\gamma)b+\gamma a\), whence
\(\gamma=-b/(a-b)\). The second identity then becomes
\(\sigma_z^2ET=(1-\gamma)b^2+\gamma a^2=-ab\).
Positive variance rules out the identically zero exceptional score.
\end{proof}
\end{annexitem}

\begin{annexitem}{annex:p15.se.oc}{thm:p15-se-oc}{The terminal exponential mean and every excess bound}
\begin{proof}
Independence and \(Ee^{hz_1}=1\) make \(e^{hZ_m}\) a
nonnegative mean-one martingale. Conditioning one step on whether the
test has already stopped proves \(Ee^{hZ_{T\wedge n}}=1\) for
every finite \(n\). On \(\{T>n\}\) its value is bounded by
\(\max(A^h,B^h)\). That part tends to zero in expectation because
termination is almost sure. The terminal contributions
\(E(e^{hZ_T};T\leq n)\) increase to the full nonnegative terminal
expectation. Subtracting the surviving part proves that this expectation
is exactly one, including its finiteness.

For \(h>0\), at an upper exit from a previous sum \(x\in(b,a)\),
put \(v=e^{-h(a-x)}\in(0,1)\). The condition is
\(e^{hz_j}\geq1/v\), and the terminal exponential normalized by
\(A^h\) has conditional mean
\(vE(e^{hz_j}\mid e^{hz_j}\geq1/v)\), between one and
\(\delta\). At a lower exit put \(v=e^{h(x-b)}>1\); its
normalized exponential has conditional mean between \(\eta\) and one.
Mix over the disjoint stopping histories. If
\(u_-=E(e^{hZ_T}\mid Q_0)\) and
\(u_+=E(e^{hZ_T}\mid Q_1)\), then
\(\eta B^h\leq u_-\leq B^h<1<A^h\leq u_+\leq\delta A^h\).
The terminal identity gives
\(\gamma=(1-u_-)/(u_+-u_-)\). This expression decreases in either
argument in their displayed ranges, giving the two bounds.
Both terminal events have positive probability, since otherwise the
terminal exponential mean would be strictly below or above one.

For \(h<0\), replace \(z\) by \(-z\), upper/lower likelihood
boundaries by \(1/B,1/A\), and the root by \(-h\). The product
\(hz\) and hence \(\eta,\delta\) are unchanged; the upper decision
probability becomes \(1-\gamma\). The positive-root bounds give exactly
the asserted negative-root formulas. Replacing both conditional means
by the boundary values gives the stated approximation. Finally
\(e^{ha}=1+ha+o(h)\) and \(e^{hb}=1+hb+o(h)\) prove its
extension at zero. That algebra does not create a missing nonzero root.
\end{proof}
\end{annexitem}

\begin{annexitem}{annex:p15.se.lattice}{thm:p15-se-lattice}{Terminal levels, complex martingales and Cramer's rule}
\begin{proof}
Before exit, the integer sum lies strictly between \(b_0\) and
\(a_0\). A downward jump has size at most \(g_1\), so a lower
exit lies between \(b_0-g_1+1\) and \(b_0\). The upper interval
follows similarly. Multiplying the root equation by \(u^{g_1}\)
gives a polynomial of degree \(g=g_1+g_2\), whose constant and leading
coefficients are the positive endpoint probabilities. Its complex roots
are therefore all nonzero, counted with multiplicity.

For each root, \(u^{Z_m}\) has conditional next expectation equal
to itself, because \(Eu^{z_1}=1\). Bounded stopping gives
\(Eu^{Z_{T\wedge n}}=1\). Every stopped sum, including a sum still
in the continuation interval, is in one fixed finite integer range.
Thus these complex variables are bounded in absolute value by one finite
constant depending on the nonzero root. The block termination argument
gives almost-sure exit, and dominated convergence yields
\(Eu^{Z_T}=1\). If \(p_j=P(Z_T=c_j)\), these equations are
\(\mathcal M p={\bf1}\). For a nonsingular matrix Cramer's rule gives the
source formula. Summing the upper terminal probabilities gives its
operating characteristic. If repeated roots give repeated rows, the
determinant is zero and that formula is not asserted.
\end{proof}
\end{annexitem}

\begin{annexitem}{annex:p15.se.finite-recursion}{thm:p15-se-finite-recursion}{A convergent absorption series and finite-prefix error certificates}
\begin{proof}
From any state in \(J\), a sufficiently long block of positive endpoint
jumps forces an upper exit. Its probability is the same positive number
at every still-continuing history. Therefore for some \(k\) and
\(c<1\), every row sum of \(K_J^{jk}\) is at most \(c^j\).
The series \(\sum_{n\geq0}K_J^n\) converges absolutely and multiplying
partial sums by \(I-K_J\) gives \(I-K_J^{n+1}\), tending to \(I\).
It is therefore the stated inverse. Conditioning on the next increment
gives the upper-exit equation \(g=v_++K_Jg\) and the expected-time
equation \(e={\bf1}+K_Je\). The inverse solves them, and the geometric
tail proves finiteness. Equivalently, their series sum the disjoint
terminal events and the survival probabilities.

By induction, \(K_J^n(x,y)\) is the probability of surviving \(n\)
steps and finishing at \(y\). Summing over \(y\) gives the survival
formula. Terminal probability vectors at the next step are obtained by
multiplying by \(v_+\) or \(v_-\). The probability generating function
for \(|s|\) in a neighborhood of one is
\[
 E_xs^T=s[(I-sK_J)^{-1}(v_++v_-)](x).
\]
The same geometric bound proves convergence there and permits termwise
derivatives of every order, giving all factorial moments and hence ordinary
moments. This argument covers a singular root matrix.

For any finite score alphabet enumerate all finite outcome words through
stage \(N\), retaining the previous continuation conditions. Sum the
probabilities of the first terminal decisions. The remaining disjoint
probability is \(P(T>N)\); any eventual decision not counted belongs to
that event. It bounds the omitted mass for either decision. The third-law
block tail, or the simple-law affinity tail, bounds this remainder before
enumeration, so no unqualified grid approximation is needed.
\end{proof}
\end{annexitem}

\begin{annexitem}{annex:p15.se.binomial}{thm:p15-se-binomial}{Two score values, strict convexity and conditional extrema}
\begin{proof}
Substitution of the two Bernoulli masses gives \(z_+,z_-\) and their
weighted mean. Conditional on \(z_1\geq u\), \(u\geq0\), the only
possible score is \(z_+\), whenever this event is nonempty. Its excess
\(z_+-u\) has supremum \(z_+\). The lower conditional calculation
gives infimum \(z_-\). For \(h>0\), the only value of \(e^{hz_1}\)
above one is \(e^{hz_+}\); taking the supremum over \(v<1\)
gives \(\delta=e^{hz_+}\). The below-one value similarly gives
\(\eta=e^{hz_-}\). Negative \(h\) reverses these two roles.

The function \(pe^{tz_+}+(1-p)e^{tz_-}\) is strictly convex,
equals one at zero, and tends to infinity at both ends. Its derivative
at zero is \(\mu\). If \(\mu\ne0\), strict convexity gives exactly
one other intersection with level one, on the side opposite that derivative.
If \(\mu=0\), zero is the unique minimum and the only intersection.
This proves the root qualification, including the midpoint exception.
\end{proof}
\end{annexitem}

\begin{annexitem}{annex:p15.se.normal}{thm:p15-se-normal}{Gaussian score, mean excess and the exponential tail-ratio extrema}
\begin{proof}
Expanding the difference between the two normal log densities gives
\(z=2dX\). The normal moment-generating integral, by completing a
square, gives \(\varphi(t)=\exp(2d\theta t+2d^2t^2)\), and
solving \(\varphi(h)=1\) gives the stated nonzero root when
\(\theta\ne0\). The mean and variance follow by scaling.

Put \(G(x)=1-\Phi(x)\) and \(H(x)=\phi(x)/G(x)\).
Integration by parts gives, for a standard normal \(Y\),
\(E(Y\mid Y\geq x)=H(x)\) and
\(E(Y^2\mid Y\geq x)=1+xH(x)\). The conditional law has a positive
density on an interval, so its variance is strictly positive. Therefore
\[
 H'(x)=H(x)(H(x)-x)<1.
\]
The mean excess \(H(x)-x\) is consequently decreasing. Conditional
on \(z\geq u\), the standardized truncation point is
\(x=u/(2d)-\theta\). Its least possible value is \(-\theta\),
so the maximum excess is \(2d\{\theta+H(-\theta)\}\).
Applying the same calculation to \(-X\) gives the negative signed
lower excess. These are exactly the stated \(\xi,\xi'\).

By reflection it suffices to take \(\theta=-\lambda<0\).
Then \(hz=2\lambda X\), and multiplication of the density of
\(N(-\lambda,1)\) by \(e^{2\lambda X}\) gives the density of
\(N(\lambda,1)\). For \(u\geq0\), the normalized upper conditional
exponential mean is
\[
 R(u)=e^{-2\lambda u}\frac{G(u-\lambda)}{G(u+\lambda)}.
\]
Its logarithmic derivative is
\(-2\lambda-H(u-\lambda)+H(u+\lambda)<0\), using the proved
bound \(H'<1\). Its supremum is therefore
\(R(0)=\Phi(\lambda)/G(\lambda)\), giving \(\delta\).
Reflection of the lower conditional expression gives \(1/R(u)\);
its infimum is \(1/\delta\). This proves the tail-ratio claims
without an unproved Mills-ratio estimate.

For completeness, in the known-variance fixed-sample comparison the
sample mean has law \(N(\theta_i,1/N)\). A level-\(\alpha\)
upper test has threshold \(\theta_0+q_{1-\alpha}/\sqrt N\).
Its alternative error is at most \(\beta\) exactly when
\(\sqrt N(\theta_1-\theta_0)\geq q_{1-\alpha}+q_{1-\beta}\).
This proves the rounded fixed-size expression in the accompanying method.
The mean log scores under the two hypotheses have magnitude
\((\theta_1-\theta_0)^2/2\). Substitution in the boundary approximation
gives its stated savings formula. These last approximations keep the
overshoot qualification; they are not exact normal SPRT sample sizes.
\end{proof}
\end{annexitem}

\begin{annexitem}{annex:p15.se.transform}{thm:p15-se-transform}{The complex martingale and its actual surviving term}
\begin{proof}
For finite \(m\), \(M_m=e^{tZ_m}\varphi(t)^{-m}\) is integrable
in absolute value by the finite real-part exponential moment and
independence. Conditioning on the next score shows
\(E(M_{m+1}\mid\mathcal F_m)=M_m\). The same one-step argument for
the stopped process gives \(EM_{T\wedge n}=1\).
Almost-sure exit excludes an identically zero score, so the earlier
block argument gives \(ET<\infty\). Put \(r=\operatorname{Re}(t)\).
If \(r\geq0\), at an upper exit
\(e^{rZ_T}\leq e^{ra}e^{rz_T^+}\), where \(z_T\) is the final
increment; at a lower exit it is at most \(e^{rb}\). By independence
from \(\{T\geq j\}\),
\[
 E e^{rz_T^+}\leq
 \sum_{j\geq1}E(e^{rz_j^+};T\geq j)
       =E e^{rz_1^+}\,ET
       \leq(1+Ee^{rz_1})ET<\infty.
\]
For \(r<0\), use the lower exit and the negative part of the final
increment in the same calculation. Thus \(Ee^{rZ_T}<\infty\).
Since \(|\varphi(t)|\geq1\), this proves the terminal integrand's
absolute integrability without a bounded-score assumption.
On \(\{T>n\}\), \(Z_n\in(b,a)\), and
\(|M_n|\leq\max(e^{\operatorname{Re}(t)b},e^{\operatorname{Re}(t)a})\)
because \(|\varphi(t)|\geq1\). This surviving expectation tends to
zero with \(P(T>n)\). If the terminal integrand is absolutely
integrable, dominated convergence on \(\{T\leq n\}\) now gives
the asserted identity. Bounded scores ensure that \(Z_T\) lies in
one fixed enlarged interval and that the terminal integrand is bounded
in absolute value, establishing the required special case. If
\(|\varphi(t)|<1\), the same decomposition works whenever the
terminal integrand is absolutely integrable and the surviving expectation
vanishes. The claimed geometric condition bounds that part by a constant
times \((\rho/|\varphi(t)|)^n\), proving the extension. For real
nonzero roots, the separate nonnegative argument in the operating-
characteristic proof establishes integrability and equality directly.

The accompanying boundary method solves a two-by-two linear system.
With \(x=\bar\gamma\widehat\Psi_+\),
\(y=(1-\bar\gamma)\widehat\Psi_-\), its equations have coefficient
rows \((A^{t_j},B^{t_j})\). Cramer's rule gives
\(x=(B^{t_2}-B^{t_1})/\mathcal D\) and
\(y=(A^{t_1}-A^{t_2})/\mathcal D\), hence the stated sum.
At zero the two roots are zero and \(h\); substitution gives
\(x=\bar\gamma\), \(y=1-\bar\gamma\), verifying conditional
normalization. The normal root expressions follow by solving
\(\mu t+\sigma_z^2t^2/2=-\tau\). This algebra proves the formula
for the declared approximation, not an error estimate for neglecting
overshoot. A transform defined with \(e^{\tau T}\) differentiates
to moments at zero on a convergence neighborhood; restriction to
\(\tau=is\) supplies the characteristic function and its factors of
\(i\). The geometric stopping tails give the moment and analytic
neighborhood for bounded scores, and for the normal examples.
For normal scores the terminal exponential moments of every real order
are finite by the final-increment argument above. Cauchy--Schwarz with
the exponential stopping tail therefore gives absolute integrability
of the terminal product for sufficiently small positive \(\tau\) as
well. The surviving geometric condition also holds near zero, completing
the required local transform extension.
\end{proof}
\end{annexitem}

\begin{annexitem}{annex:p15.se.continuous-density}{thm:p15-se-continuous-density}{An explicit integral and the approximation's moments}
\begin{proof}
For \(s>0\), let
\(F(s)=\int_0^\infty e^{-v^2-s^2/v^2}\,dv\).
Differentiation under the integral is justified on every compact
positive \(s\)-interval by the exponential decay at zero and infinity.
Substitution \(w=s/v\) gives
\(\int_0^\infty v^{-2}e^{-v^2-s^2/v^2}\,dv=F(s)/s\),
so \(F'(s)=-2F(s)\). Dominated convergence at zero gives
\(F(0)=\sqrt\pi/2\), hence \(F(s)=(\sqrt\pi/2)e^{-2s}\).
Putting \(v=\sqrt{a_*/m}\), for \(a_*,b_*>0\), yields
\[
 \int_0^\infty m^{-3/2}e^{-a_*/m-b_*m}\,dm
          =\frac{\sqrt\pi}{\sqrt{a_*}}\,e^{-2\sqrt{a_*b_*}}.
\]
Apply this with \(a_*=c^2/4\), \(b_*=1+q\), and multiply by
\(ce^c/(2\sqrt\pi)\). This gives the claimed Laplace transform.
At \(q=0\) it is one, proving normalization. The density's exponential
decay at both endpoints permits its first two moment integrations and
differentiation at zero. The first derivative of its log Laplace transform
is \(-c/(2\sqrt{1+q})\), and the second is
\(c/(4(1+q)^{3/2})\). At zero these give mean \(c/2\) and
variance \(c/4\). No statement about the error of the discrete
stopping-time approximation follows from normalization alone.
\end{proof}
\end{annexitem}

\begin{annexitem}{annex:p15.se.deadline}{thm:p15-se-deadline}{Endpoint events, first exits and exact normal marginals}
\begin{proof}
If \(Z_N\leq b\), some boundary must have been crossed by stage
\(N\): otherwise every score sum through that stage would still lie
in \((b,a)\). This proves the null inequality; the upper event gives
the alternative inequality by the same argument. It is not necessary
that the first exit have the same direction as the score at stage
\(N\). Under normal increments, \(Z_N\) has the exact normal law
with mean \(N\mu_i\) and variance \(N\sigma_i^2\). Standardization
gives the two expressions. Normality is used only for that calculation,
not for the event inclusions.
\end{proof}
\end{annexitem}

\begin{annexitem}{annex:p15.se.truncation}{thm:p15-se-truncation}{Which errors can be introduced at the cap}
\begin{proof}
Couple the capped and uncapped rules on the same infinite observation
sequence. Their decisions agree if the boundary test stops by stage
\(N\). A new null error can therefore occur only when the test is
still continuing after stage \(N\), \(0<Z_N<a\), and the uncapped
test would later choose zero. Its probability is no larger than
\(P_0(0<Z_N<a)\). Any capped null rejection otherwise already belongs
to the uncapped rejection event. Taking the union bound gives the null
inequality. The new alternative error is treated with
\(b<Z_N\leq0\) and a later uncapped upper decision. This gives the
second inequality and includes the zero-score tie rule.
The normal interval probabilities follow by the same exact standardization
as in the deadline proof. The uncapped base errors are actual decision
probabilities, so replacing them by nominal targets without a separate
guarantee changes the claimed bound. The exact upper substitutes follow
from the earlier coupled-error theorem.
\end{proof}
\end{annexitem}

\begin{annexitem}{annex:p15.se.repeated-fixed}{thm:p15-se-repeated-fixed}{Every look is valid, but their first rejection has a different level}
\begin{proof}
At a preselected stage \(j\), uniformity gives
\(P(U_j\leq\alpha)=\alpha\). The rule fails to reject through
\(M\) stages exactly when each of the independent variables exceeds
\(\alpha\), with probability \((1-\alpha)^M\). Complementing gives
the formula. At the stated values it is
\(1-(19/20)^{20}\approx.6415140776\). As \(M\to\infty\),
the no-rejection probability tends to zero, so unbounded inspection
rejects almost surely. This first-crossing rule is measurable from the
observed history, but its error is the union probability, not the error
at any one look. The SPRT comparison is the already proved bound on
the entire upper-terminal event with its chosen conservative boundary.
No independence claim is made for overlapping cumulative samples.
\end{proof}
\end{annexitem}

\begin{annexitem}{annex:p15.se.information}{thm:p15-se-information}{The original conditional-density comparison and its unbounded extension}
\begin{proof}
First suppose a competing test stops by a finite horizon \(N\).
The stopped-history density ratio is \(L_T\), since the common
decision/stopping region multiplies both history densities. Summing over
the finite terminal stages gives
\[
 E_0(L_T;D=0)=v,\quad E_0(L_T;D=1)=1-v,
\]
and, by reversing the likelihood,
\[
 E_1(L_T^{-1};D=0)=1-u,\quad
 E_1(L_T^{-1};D=1)=u.
\]
An independent seed has the same density under both laws and cancels.
For any positive integrable variable \(Y\) with integrable log,
\(E\log Y\leq\log EY\): put \(r=EY>0\), use
\(\log(Y/r)\leq Y/r-1\), and integrate. Apply this inequality
on each terminal event. Under \(P_0\) it gives
\[
 E_0Z_T\leq(1-u)\log\frac v{1-u}
                    +u\log\frac{1-v}u.
\]
Under \(P_1\), applying it to the reversed likelihood gives
\[
 E_1(-Z_T)\leq v\log\frac{1-u}v
                    +(1-v)\log\frac u{1-v}.
\]
Bounded-time predictability and integrable scores give
\(E_iZ_T=\mu_iE_iT\). Dividing by the negative or positive
mean proves both finite-horizon inequalities with the asserted signs.

For an arbitrary almost-surely terminating competing test, cap its
stopping time at \(N\), keeping its decision if already terminal and
choosing either measurable decision on the remaining histories. The two
errors \(u_N,v_N\) converge to \(u,v\), because their disagreement
probability is at most \(P_i(T>N)\), tending to zero. Its expected
time is \(E_i(T\wedge N)\leq E_iT\). Apply the finite-horizon
inequalities and let \(N\to\infty\). The right sides are continuous
at the assumed strictly interior errors, and monotone convergence gives
the expected-time limit, possibly infinite. This fills the source's
unbounded stopping extension without assuming a terminal entropy identity
for an unintegrable competing procedure.
\end{proof}
\end{annexitem}

\begin{annexitem}{annex:p15.se.no-overshoot}{thm:p15-se-no-overshoot}{Exact attainment and a bounded conditional Jensen gap}
\begin{proof}
Under the stated simple-law assumptions the SPRT has finite expected
termination, so its stopped-sum identity applies. If it has no overshoot,
the density comparisons become equalities:
\(1-\beta'=A\alpha'\) and
\(\beta'=B(1-\alpha')\). Thus the two possible values of \(Z_T\)
are exactly the logarithms appearing in the conditional information
calculation, and the inequalities there are equalities. The SPRT attains
both universal lower bounds at its actual errors. No test with those
same errors can improve either expected size. A walk with jumps
\(d,-d\) cannot cross an aligned boundary without landing on it,
which proves the displayed example.

For bounded increments, on each terminal decision the value \(Z_T\)
lies in an interval of width at most \(C\): \([b-C,b]\) or
\([a,a+C]\). For any real random variable \(Y\) in an interval
\([v,v+C]\),
\[
        0\leq\log Ee^Y-EY\leq C.
\]
The lower bound is the logarithmic expectation inequality; the upper
bound follows from \(\log Ee^Y\leq v+C\) and \(EY\geq v\).
Apply this on each decision event with \(Y=Z_T\) under the null,
and \(Y=-Z_T\) under the alternative. The conditional density
identities make the weighted logarithmic terms exactly the binary
information numerators. Their weighted Jensen gaps are therefore
\(|\mu_i|E_iT-I_i\), between zero and \(C\). Division proves
the quantitative comparison. The target is the information lower bound,
so no unproved general exact-optimality assertion is introduced.
\end{proof}
\end{annexitem}
